\documentclass[11pt]{article}
\usepackage[a4paper,margin=25mm]{geometry}
\usepackage{amsmath,amssymb,amsthm,hyperref}
\emergencystretch=2em
\newtheorem{theorem}{Theorem}\newtheorem{lemma}{Lemma}
\newcommand{\E}{\mathbb E}
\title{Endpoint profiles on their natural scale}
\author{Complete theorem; independently reviewed}
\date{30 September 2026}
\begin{document}\maketitle

\begin{theorem}\label{main}
Under the monotone mixed-moment hypotheses of
\path{optimal_interior_endpoint_profile.tex}, put
$u_i=1-p_i$ and $s=\sum_i u_i$. For every fixed $L<\infty$,
uniformly over rows with $0\le s\le L$,
\[
 \max_i|m u_i-s|\le C_L\left(\sqrt{\frac{s}{m}}+\frac1m\right).
\]
\end{theorem}

The two inputs beyond the reviewed fixed-annulus profile theorem are
the sharp relative balance and endpoint distribution bound in
\path{../sharp_endpoint_relative_balance.tex}, and the variable-center
Jackson estimates in
\path{../rational_designs/asymptotic/variable_jackson_endpoint_kernel.tex}.
The present note proves the measure and monotonicity assembly.

\paragraph{Coordinates and sharp bounds.}
Retain $M=m-1$, $\ell=\lfloor M/4\rfloor$, $N=M-\ell$,
$z_i=\ell u_i$, and the marked row probabilities
\[
 x_i=M^{-1}\sum_{j\ne i}z_j,\quad h_i=z_i-x_i,\qquad
 d\mu_i=(\ell+1)\E_B\prod_{j\in B}p_j\,dw.
\]
Here $B$ ranges over the $\ell$-subsets of the unmarked indices.
The reviewed sharp bounds give, uniformly in $i$,
\begin{equation}\label{sharp}
 \max_j z_j\le C(x_i+m^{-1}),\qquad
 \mu_i([0,t])\le C(t+m^{-1})\quad(t\ge0).
\end{equation}
For the first bound, absorb the marked column's contribution in
the full-family balance estimate. For the second apply the sharp
CDF lemma to the deleted $M$-column tensor and note that
$\mu_i\le C M w$, with only a fixed rescaling of its deficit
coordinate. Also $\mu_i$ has total mass one.

For centers $a\ge T/m$, with $T$ a sufficiently large fixed constant,
write
\[
 b_a=\sqrt{a/m},\qquad
 E_a(x)=b_a^{-1}(1+|x-a|/b_a)^{-2r}
 \quad\text{when }x\in[c a,Ca].
\]
Fixed changes in the constants or scale of this envelope are harmless.
The Jackson degree is a sufficiently small fixed multiple of $m$,
and its fixed power $r>4$ can be increased for the negligible global
high-order remainder. Let $Q_a$ be the corresponding positive
polynomial, normalized so its scalar beta-reference integral lies
between two positive absolute constants. Its near-diagonal lower
bound is $Q_a\ge c/b_a$ on $|x-a|\le c b_a$.

\paragraph{The variable localized comparisons.}
For $T/m\le a\le a_*$, where $a_*>0$ is fixed and sufficiently
small, the variable-kernel note gives, with $H=\max_j z_j-\min_jz_j$,
\begin{align}
 |\mu_i(Q_a)-\beta_\ell(Q_a)|
 &\le \frac C a\int_{x_i\in[c a,Ca]}H^2E_a(x_i)\,d\mu_i
       +C(ma)^{1/2-r}+m^{-A},\label{unmarked}\\
 |\mu_i(h_iQ_a)|
 &\le C\int_{x_i\in[c a,Ca]}
 \left(\frac{H^2}{\sqrt{ma}}+\frac{H^3}{a}\right)
 E_a(x_i)\,d\mu_i
       +Ca(ma)^{1/2-r}+m^{-A}.
 \label{marked}
\end{align}
Here $A>10$ is fixed, and the constants are independent of the
varying center. The same relative tails bound scalar kernel integrals
outside the displayed relative annulus, including an extra factor
$|h_i|$ in the marked case. These follow from \eqref{sharp}, not from
the total-mass bound alone. On the relative annulus,
\eqref{sharp} gives $H\le Ca$ because $ma\ge T$.

\begin{lemma}[Uniform variable-width local mass]\label{mass}
After choosing $T$ sufficiently large and then a sufficiently small
fixed $a_0>0$, uniformly for $T/m\le a\le a_0$ and all marked columns,
\[
 \mu_i(|x_i-a|\le b_a)\le C b_a,
 \qquad c\le\mu_i(Q_a)\le C,
 \qquad\int_{x_i\in[c a,Ca]}E_a(x_i)\,d\mu_i\le C.
\]
Moreover, the portion of the last integral outside
$|x_i-a|\le B b_a$ is at most $C B^{1-2r}$, up to the separately
controlled relative-annulus tails.
\end{lemma}
\begin{proof}
First prove only the upper bounds. Choose a fixed small $a_1$ for
which the coefficient $Ca$ in \eqref{unmarked} is small enough for
absorption on $a\le a_1$. Consider the supremum
\[
 V=\sup_{i,\,T/m\le a\le a_1}
       \mu_i([a-c b_a,a+c b_a])/b_a.
\]
For centers below $T/m$, \eqref{sharp} bounds this normalized local
mass by a fixed constant depending on $T$, using width
$\sqrt{a/m}+1/m$. For centers in a fixed interval above $a_1$, the
already proved fixed-annulus mass bound supplies a finite uniform
constant. Thus, for any center in the supremum, intervals used to
cover $[ca,Ca]$ have normalized local mass at most $V+C_T$; here
$C_T$ can also depend on the fixed upper boundary $a_1$.
Widths on this relative annulus are comparable to $b_a$.
Summing the polynomial envelope over these intervals gives
\[
 \int_{[ca,Ca]}E_a\,d\mu_i\le C(V+C_T).
\]
The lower bound for $Q_a$ on its central window, its bounded
reference integral, and \eqref{unmarked} now imply
\[
 V\le C+C a_1(V+C_T)+C T^{1/2-r}.
\]
Absorb $C a_1 V$. This proves a finite upper bound $V\le C_1$,
uniform in $m$ and the marked column. It gives the stated upper
interval and envelope bounds for all centers below $a_1$.

Now choose $T$ larger, if necessary, so the relative tail term in
\eqref{unmarked} is smaller than a fixed fraction of the positive
reference mass. The preceding upper bound remains finite with this
choice. Choose $a_0\le a_1$ still smaller so that $C a_0 C_1$ is
also smaller than that fraction. Equation \eqref{unmarked} proves
the lower kernel-mass bound for $T/m\le a\le a_0$.
The tail estimate follows by summing the local interval bounds over
length-$b_a$ intervals. This two-stage choice avoids assuming that
fixed-annulus constants remain bounded when its lower endpoint
tends to zero.
\end{proof}

\paragraph{A normalized maximum.}
Use the common coordinate and deviations
\[
 y=\ell s/m,\quad h_i^*=z_i-y,\quad
 x_i=y-h_i^*/M,\quad h_i=(m/M)h_i^*.
\]
Put $b(y)=\sqrt{y/m}+1/m$. Sharp balance gives
\begin{equation}\label{hbound}
 |h_i^*|\le C(y+1/m),\qquad
 |x_i-y|\le C(y+1/m)/m=o(b(y))
\end{equation}
uniformly on bounded $y$-intervals. The same bounds hold for any row
entering the corresponding deleted-coordinate interval, by absorbing
the marked contribution as in \eqref{sharp}.

Fix a bounded $y$-interval containing the desired range and a fixed
enlargement. Consider the maximum
\[
 D=\max_{i,\,\text{rows in that interval}} |h_i^*(y)|/b(y).
\]
Rows with $my$ bounded have bounded normalized deviation by
\eqref{hbound}. Rows with $y$ bounded below by any fixed positive
constant have bounded normalized deviation by the already proved
interior square-root theorem. Consequently, if $D$ exceeds a
sufficiently large fixed constant, a maximizing row $y_0$ lies in
\[
 2T/m\le y_0\le a_2,
\]
where $a_2>0$ can be chosen as small as needed after Lemma~\ref{mass}.
Write $\Delta=|h_i^*(y_0)|$ and $b_0=b(y_0)$. Then
$\Delta=D b_0$ and $\Delta\le Cy_0$.
For rows whose coordinates lie in a fixed relative neighborhood of
$y_0$, maximality gives
\begin{equation}\label{range}
 |h_j^*|\le C\Delta\quad\text{for every }j,
 \qquad H\le C\Delta.
\end{equation}
This also holds in the relative neighborhood of the selected
deleted coordinate $x_0$, by \eqref{hbound}.

\paragraph{A favorable one-sided window.}
Choose fixed $B_0$ sufficiently large and $A_0=2B_0$.
Center the test at $a=x_0+A_0b_0$ for a positive selected deviation,
and at $a=x_0-A_0b_0$ for a negative one. Since $\Delta=D b_0\le Cy_0$, one has $y_0/b_0\ge D/C$.
Choosing the fixed contradiction threshold for $D$ after $B_0$ therefore
keeps these centers and windows in a fixed relative neighborhood of
$y_0$, without changing the already fixed $T$. Then $a\asymp y_0$ and
$b_a\asymp b_0$. The favorable window lies strictly on the selected
side, so monotonicity gives the chosen sign of $h_i$ with magnitude
at least $\Delta/2$ there, whenever $D$ is sufficiently large.

Lemma~\ref{mass} and its local tails imply a positive kernel mass
on that favorable window. Within the relative annulus,
\eqref{range} bounds the opposite-sign contribution outside it by
$C\Delta B_0^{1-2r}$, made small by choosing $B_0$ large.
Outside the relative annulus the scalar marked tail is at most
\[
 Ca(ma)^{1/2-r}=C b_a(ma)^{1-r}.
\]
Since $ma\ge T$ and $\Delta\asymp D b_a$, this is a small fraction
of $\Delta$ after the fixed choices. Hence
\begin{equation}\label{positive}
 |\mu_i(h_iQ_a)|\ge c\Delta.
\end{equation}

The opposite bound \eqref{marked}, the local envelope mass estimate,
and \eqref{range} give
\[
 |\mu_i(h_iQ_a)|
 \le C\left(\frac{\Delta^2}{\sqrt{ma}}+
                         \frac{\Delta^3}{a}\right)
       +C b_a(ma)^{1-r}+m^{-A}.
\]
Since $\Delta\le Ca$ and $a\le C a_2$, the ratio of the first part
to $\Delta$ is at most $C\sqrt{a/m}+Ca$.
Choose $a_2$ sufficiently small and then $m$ large. The tail is
absorbed using $ma\ge T$ and the choice of $D$ above the fixed
threshold. Finally $\Delta\ge c/m$, so $m^{-A}$ is negligible.
This contradicts \eqref{positive}. Therefore $D$ is bounded by a
constant depending only on the desired upper endpoint $L$.
The relation $m u_i-s=(m/\ell)h_i^*$ and $y\asymp s$ proves
Theorem~\ref{main}.
\paragraph{Relative synchronization and the opposite endpoint.}
For $s>0$, the theorem gives
\[
 \max_i\left|\frac{m u_i}{s}-1\right|
 \le C_L\left((ms)^{-1/2}+(ms)^{-1}\right).
\]
Thus the scaled columns synchronize relatively whenever $ms\to\infty$
and $s\le L$. Reversing the row order and replacing every $p_i$ by
$1-p_i$ preserves the mixed moments, since
$\int_0^1(1-t)^r\,dt=1/(r+1)$.
The same conclusion therefore holds at the opposite endpoint, with
$s$ replaced by $\sum_i p_i$ and $u_i$ by $p_i$.

The two orders in the theorem are necessary in this generality.
The reviewed notes \path{../endpoint_profile_fluctuation_barrier.tex}
and \path{../finite_profile_realization.tex} give fluctuations of order
$m^{-1/2}$ at fixed positive $s$. The note
\path{../endpoint_profile_sharpness.tex} gives actual finite uniform
dice with $s\asymp m^{-2}$ and error $\asymp m^{-1}$, so the additive
floor cannot be removed. The subsequent note
\path{intermediate_endpoint_sharpness.tex} gives matching finite examples
at every positive bounded scale for the refined envelope
$\min\{ms,\sqrt{s/m}+1/m\}$, using the trivial upper bound $ms$.
The realized deficit is comparable to the prescribed scale; no exact
prescribed numerical value of $s$ is asserted.
\end{document}
